\section{Список аниме/манги/VN для просмотра именно как материала по архетипам}

Список собран из всех тайтлов, к которым книга обращалась по ходу глав 0--11 — не рейтинг "лучшего аниме", а рабочий материал, каждый пункт которого иллюстрирует конкретный, названный механизм.

\textbf{История и теория (главы 0--1):} Astro Boy / Tetsuwan Atom; Mazinger Z; Gundam; Sally the Witch; Sailor Moon; Creamy Mami; Macross; Neon Genesis Evangelion; Lucky Star; K-On!.

\textbf{Dere-семейство (глава 3):} Cardcaptor Sakura; Toradora!; Naruto; Mirai Nikki / Future Diary; Zero no Tsukaima; Dragon Ball.

\textbf{Амплуа (глава 4):} Fruits Basket; Kaguya-sama: Love Is War; Revolutionary Girl Utena; Oreimo; Haruhi Suzumiya (сериал).

\textbf{Оси характера (глава 5):} Bakemonogatari (Monogatari series); Attack on Titan; Steins;Gate; Violet Evergarden; Ghost in the Shell; Sousou no Frieren.

\textbf{Визуальный и голосовой код (главы 6--7):} Spice and Wolf; Kiki's Delivery Service; Fullmetal Alchemist; Spirited Away.

\textbf{Близость (глава 8):} Bloom Into You / Yagate Kimi ni Naru; Mushishi; Given; Oshi no Ko; Chobits; Ijiranaide, Nagatoro-san; Barakamon.

\textbf{Тёмные архетипы (глава 9):} Higurashi no Naku Koro ni; Wonder Egg Priority; Vivy: Fluorite Eye's Song; Re:Zero; Plastic Memories.

\textbf{Современные гибриды (11.12):} Fate (серия); Spy x Family; My Dress-Up Darling; Bocchi the Rock!; Chainsaw Man; Kusuriya no Hitorigoto / The Apothecary Diaries.

Сознательный уклон в сторону не-сёнэн тайтлов и разнообразия жанров — не эстетическая прихоть, а прямое следствие принципа этой книги (0.5): точность соответствия важнее жанровой мейнстримности.
